\chapter{Desarrollo}
\label{chap:desarrollo}

Este capítulo constituye el núcleo técnico de la memoria.
En él se describe, de manera pormenorizada, el proceso de construcción del sistema de navegación volumétrica siguiendo la metodología iterativa e incremental expuesta en la Sección~\ref{sec:metodologiadesarrollo}.
Cada una de las secciones que siguen se corresponde con una iteración del desarrollo y se articula en torno a un mismo esquema: el objetivo perseguido, las decisiones de diseño junto con su implementación, las particularidades técnicas más destacables y la validación del incremento resultante.

El recorrido se inicia con una iteración preparatoria que asienta la arquitectura modular y el entorno de trabajo (Iteración~0).
Sobre ella se erige un producto mínimo viable, constituido por la representación y la construcción del espacio navegable junto con una primera versión funcional de la búsqueda de rutas (Iteraciones~1 a~3), que se amplía después de forma progresiva hasta alcanzar un producto completo y optimizado.
En coherencia con la planificación descrita en la Sección~\ref{sec:planificacioninicial}, cada una de las fases de implementación allí previstas se materializa aquí a través de una o varias iteraciones.
Así, una vez alcanzado el núcleo mínimo, las iteraciones restantes abordan sucesivamente el postprocesado de trayectorias, la persistencia en disco, la robustez geométrica, la integración en el motor, la optimización del rendimiento y, finalmente, la evasión local entre agentes.
El capítulo concluye con una discusión global que sintetiza la arquitectura resultante y las principales decisiones de diseño.

\section{Iteración 0: Cimientos del proyecto}

La primera iteración no produce funcionalidad de navegación, sino que establece los cimientos sobre los que se sustentará todo el desarrollo posterior.
Su objetivo es doble: definir una arquitectura modular que favorezca el desacoplamiento y poner en marcha el flujo de trabajo y las garantías de calidad descritas en el capítulo anterior.

El sistema se organiza en módulos con responsabilidades bien delimitadas, separados físicamente mediante \textit{assemblies} (las unidades de compilación de C\#).
Esta separación no es meramente organizativa: al aislar el código de edición y de pruebas en \textit{assemblies} independientes, se garantiza que estos queden excluidos de las compilaciones finales del producto, sin penalizar el tamaño ni el rendimiento del ejecutable distribuido.
La estructura resultante se divide en tres grandes bloques, ilustrados en la Figura~\ref{fig:arquitectura}:

\begin{itemize}
    \item \textbf{Runtime:} contiene la totalidad de la lógica de navegación (representación del espacio, construcción y búsqueda de rutas). Es el único módulo que se incluye en la versión final del producto.
    \item \textbf{Editor:} alberga las herramientas de configuración y visualización integradas en el editor de Unity. Depende del módulo \textit{Runtime}, pero nunca a la inversa.
    \item \textbf{Tests:} agrupa las pruebas automatizadas, divididas a su vez en pruebas en modo de edición (\textit{EditMode}) y en modo de ejecución (\textit{PlayMode}).
\end{itemize}

Internamente, el módulo \textit{Runtime} se subdivide a su vez en espacios de nombres cohesivos: \texttt{Core} (las estructuras de datos fundamentales), \texttt{Builder} (la construcción del volumen), \texttt{Pathfinding} (la búsqueda y el postprocesado) y \texttt{Unity} (los componentes que exponen el sistema al motor).
Esta organización inicial se ampliaría más adelante, a medida que el sistema crecía, con nuevos espacios de nombres —en particular \texttt{Avoidance} (la evasión local), incorporado en la última iteración—; la arquitectura definitiva se recoge en la discusión final (Figura~\ref{fig:arquitectura-final}).

\begin{figure}[htp]
  \centering
  \begin{tikzpicture}[
      font=\small,
      modulo/.style={rounded corners=3pt, draw, thick, align=center,
                     inner sep=8pt, minimum width=3.4cm, minimum height=1.5cm},
      contenedor/.style={rounded corners=4pt, draw=udcpink, thick, fill=udcpink!10,
                         inner sep=10pt},
      ns/.style={rounded corners=2pt, draw=udcpink!60!black, fill=white, align=center,
                 minimum width=2.1cm, minimum height=9mm, font=\footnotesize\ttfamily},
      dep/.style={dashed, -{Stealth[length=3mm]}, thick, draw=black!70}
    ]

    % Espacios de nombres internos del modulo Runtime
    \node[ns] (core) {Core};
    \node[ns, right=3mm of core]    (builder) {Builder};
    \node[ns, right=3mm of builder] (path)    {Pathfinding};
    \node[ns, right=3mm of path]    (unity)   {Unity};

    % Contenedor Runtime (en capa de fondo para no tapar los nodos)
    \begin{scope}[on background layer]
      \node[contenedor, fit=(core)(builder)(path)(unity),
            label={[font=\bfseries]above:Runtime}] (runtime) {};
    \end{scope}

    % Modulos Editor y Tests
    \node[modulo, fill=ficblue!15, draw=ficblue,
          above=16mm of runtime.north west, anchor=south west]
      (editor) {\textbf{Editor}\\[2pt]\footnotesize Herramientas de\\configuración y\\visualización};
    \node[modulo, fill=udcgray!40, draw=udcgray!80!black,
          above=16mm of runtime.north east, anchor=south east]
      (tests) {\textbf{Tests}\\[2pt]\footnotesize Pruebas \textit{EditMode}\\y \textit{PlayMode}};

    % Dependencias unidireccionales hacia Runtime
    \draw[dep] (editor.south) -- (editor.south |- runtime.north);
    \draw[dep] (tests.south)  -- (tests.south  |- runtime.north);

  \end{tikzpicture}
  \caption[Arquitectura modular del sistema y sus dependencias.]
  {Arquitectura modular del sistema y sus dependencias.
  Los módulos \textit{Editor} y \textit{Tests} dependen del módulo \textit{Runtime} (flechas discontinuas), pero nunca a la inversa;
  únicamente \textit{Runtime} se incluye en la versión final del producto.}
  \label{fig:arquitectura}
\end{figure}

Sobre estos cimientos se configuró el flujo de trabajo basado en \textit{pull requests} y el sistema de \acrshort{ci}, descritos en la Sección~\ref{sec:metodologiadesarrollo}.
Una de las primeras garantías de calidad que ejecuta este sistema de \acrshort{ci} es la verificación automática del estilo: ante cada \textit{pull request}, una comprobación de formato (\textit{format check}) con el formateador CSharpier impide la integración del código si este no se ajusta a la guía de estilo común adoptada.
De este modo, la deuda técnica asociada al formato y la consistencia del código se mantuvo controlada, sin intervención manual, a lo largo de todo el proyecto.

\section{Iteración 1: Representación del espacio navegable}
\label{sec:iter-representacion}

El objetivo de esta iteración es implementar la estructura de datos que representa el espacio navegable de forma compacta.
Tal y como se justificó en la Sección~\ref{sec:representacionespacionavegable}, la técnica seleccionada es el \acrshort{svo}, una variante del \gls{octree} que solo subdivide las regiones del espacio que contienen geometría.

\subsection{Organización jerárquica}

El \acrshort{svo} se modela como un conjunto de capas (\textit{layers}), donde cada capa es un \textbf{vector plano} de nodos ordenado por su código de Morton.
Las capas superiores representan nodos de gran tamaño (espacio mayoritariamente libre), mientras que las inferiores aumentan progresivamente la resolución.
La capa más baja la constituyen las \textbf{hojas}, que no almacenan ocho hijos como los nodos internos, sino una rejilla de $4\times4\times4$ \glspl{voxel} (véase la Figura~\ref{fig:hoja-mascara}).

Esta decisión de diseño es deliberada y responde a los principios del \acrshort{dod} expuestos en la Sección~\ref{sec:dots}.
Almacenar cada hoja como una rejilla densa de 64 \glspl{voxel}, en lugar de seguir subdividiendo el árbol cuatro niveles más, reduce drásticamente el número de nodos y, sobre todo, mantiene los datos contiguos en memoria.
La búsqueda de un nodo dentro de una capa se resuelve mediante una \textbf{búsqueda binaria} sobre el vector ordenado, aprovechando que los códigos de Morton preservan la localidad espacial.

\begin{figure}[htp]
  \centering
  \begin{tikzpicture}[font=\small]
    % Una capa 4x4 de la hoja; voxeles ocupados en rosa
    \foreach \c/\r in {1/0, 3/0, 0/2, 2/2, 2/3}
      \fill[ocupado] (\c*0.72,-\r*0.72) rectangle ++(0.72,-0.72);
    \foreach \c in {0,1,2,3}
      \foreach \r in {0,1,2,3}
        \draw[celda] (\c*0.72,-\r*0.72) rectangle ++(0.72,-0.72);
    \node[font=\footnotesize, anchor=south] at (1.44,0.12) {Hoja (una capa $4\times4$)};

    % Flecha de codificacion
    \draw[flecha] (3.2,-1.44) -- ++(1.45,0) node[midway, above, font=\scriptsize] {1 bit/vóxel};

    % Mascara de 64 bits (8x8)
    \begin{scope}[xshift=5cm, yshift=0.35cm]
      \foreach \c/\r in {1/0,3/0,5/1,0/2,2/2,6/3,2/3,4/4,7/5,1/6,3/6,5/7}
        \fill[ocupado] (\c*0.34,-\r*0.34) rectangle ++(0.34,-0.34);
      \foreach \c in {0,...,7}
        \foreach \r in {0,...,7}
          \draw[celda] (\c*0.34,-\r*0.34) rectangle ++(0.34,-0.34);
      \node[font=\footnotesize, anchor=south] at (1.36,0.12) {Máscara de 64 bits (\texttt{ulong})};
    \end{scope}
  \end{tikzpicture}
  \caption[Hoja del \acrshort{svo} y su máscara de ocupación de 64 bits.]
  {Hoja del \acrshort{svo} y su máscara de ocupación de 64 bits.
  Cada hoja apila cuatro capas de $4\times4$ vóxeles; cada uno de esos $4\times4\times4=64$ vóxeles se corresponde con un bit de un único entero de 64 bits (\texttt{ulong}): los ocupados (en rosa) fijan su bit a~1 y los libres a~0.}
  \label{fig:hoja-mascara}
\end{figure}

\subsection{Codificación de las hojas mediante manipulación de bits}

Cada hoja almacena su ocupación en un \textbf{único entero de 64 bits} (\texttt{ulong}), donde cada bit indica si el \gls{voxel} correspondiente está libre u ocupado.
Esta codificación, que materializa la técnica de manipulación de bits introducida en los fundamentos, permite que una hoja completa ocupe tan solo 8 bytes y que operaciones globales como ``¿está la hoja completamente vacía?'' o ``¿está completamente llena?'' se resuelvan con una única comparación entera, sin necesidad de inspeccionar los 64 \glspl{voxel} de forma individual.

La conversión entre las coordenadas tridimensionales de un \gls{voxel} dentro de la hoja y su índice de bit se realiza también mediante desplazamientos, empaquetando las tres coordenadas (de 2 bits cada una) en los 6 bits del índice.

\subsection{Códigos de Morton}

Para mapear las coordenadas tridimensionales $(x, y, z)$ de un nodo a un índice escalar que preserve la localidad espacial se emplean los códigos de Morton, basados en las curvas de orden Z descritas en la Sección~\ref{sec:representacionespacionavegable} (véase la Figura~\ref{fig:morton_curve}).
La codificación se obtiene intercalando los bits de las tres coordenadas mediante una secuencia de operaciones de desplazamiento y máscara, tal y como se muestra en la Figura~\ref{fig:morton-interleave}.

\begin{figure}[H]
\begin{lstlisting}[language={[Sharp]C}]
// Intercala dos ceros entre cada bit del valor de entrada.
static uint Interleave00(uint x)
{
    x = (x ^ (x << 16)) & 0xFF0000FF;
    x = (x ^ (x <<  8)) & 0x0300F00F;
    x = (x ^ (x <<  4)) & 0x030C30C3;
    x = (x ^ (x <<  2)) & 0x09249249;
    return x;
}

// El codigo final entrelaza los bits de las tres coordenadas.
public MortonCode(int x, int y, int z) =>
    _code = Interleave00((uint)x)
          | (Interleave00((uint)y) << 1)
          | (Interleave00((uint)z) << 2);
\end{lstlisting}
\caption{Codificación de un código de Morton mediante intercalado de bits.}
\label{fig:morton-interleave}
\end{figure}

Sobre esta representación, operaciones jerárquicas habituales se vuelven triviales: el código del nodo padre se obtiene desplazando cada coordenada un bit a la derecha, y el de un hijo concreto, desplazándola a la izquierda e incorporando el bit correspondiente al octante.
Dado que la codificación emplea enteros de 32 bits, la resolución máxima es de $1024$ ($2^{10}$) celdas por eje, lo que acota el número máximo de capas del árbol.

\subsection{Punteros compactos: el enlace del SVO}
\label{sec:svolink}

Un último elemento fundamental de esta iteración es el mecanismo de referencia entre nodos.
En lugar de emplear referencias de C\# (punteros de 8 bytes que dispersarían los datos en la \gls{heap}), se diseñó un \textbf{puntero compacto} de 32 bits que codifica, en un único entero, toda la información necesaria para localizar cualquier nodo o \gls{voxel} del árbol.
El reparto de bits se muestra en la Figura~\ref{fig:svolink-bits}.

\begin{figure}[H]
\begin{lstlisting}[language={[Sharp]C}]
const uint _IS_NODE_MASK       = 1u << 31;       // bit 31:    tipo (1 nodo, 0 voxel)
const uint _OFFSET_MASK        = 0x1FFFFFF << 6; // bits 30..6: offset (25 bits)
const uint _CONCRETE_DATA_MASK = 0x3F;           // bits  5..0: capa o subnodo (6 bits)
\end{lstlisting}
\caption{Disposición de los 32 bits del enlace (\texttt{SVOLink}).}
\label{fig:svolink-bits}
\end{figure}

El bit más significativo distingue si el enlace apunta a un nodo interno o a un \gls{voxel}; los 25 bits centrales almacenan el desplazamiento dentro de la capa o del vector de hojas; y los 6 bits inferiores reutilizan su significado según el tipo: la capa del nodo o el índice del subnodo dentro de la hoja.
Esta compacidad no solo ahorra memoria, sino que permite que el enlace se utilice directamente como clave en las colecciones \textit{hash} del buscador de rutas.
Para que esto resulte eficiente, el tipo implementa la interfaz \texttt{IEquatable}, evitando así el \textit{boxing} y la comparación por reflexión en la ruta más caliente del algoritmo.

Una vez introducidos todos sus elementos, la Figura~\ref{fig:core-clases} sintetiza las estructuras de datos del espacio de nombres \texttt{Core} y las relaciones de composición y de referencia que las vinculan.

\begin{figure}[htp]
  \centering
  \begin{tikzpicture}[
      font=\small,
      clase/.style={draw=ficblue!75!black, fill=ficblue!7, thick, rounded corners=1pt,
                    rectangle split, rectangle split parts=3, align=left,
                    text width=3.3cm, inner xsep=5pt, inner ysep=3.5pt, font=\scriptsize},
      comp/.style={draw=ficblue!75!black, thick,
                   -{Diamond[length=2.8mm, width=1.9mm, fill=ficblue!75!black]}},
      deref/.style={draw=black!65, dashed, -{Stealth[length=2.2mm]}, thick},
      etiq/.style={font=\scriptsize, fill=white, inner sep=1pt},
      mult/.style={font=\scriptsize, inner sep=1.5pt}
    ]

    % --- Clases ---
    \node[clase] (svo) at (0,0) {
      {\itshape class}\\[-1pt]\textbf{SVO}
      \nodepart{two}\ttfamily Layers : SVONode[][]\\ LeafNodes : SVOLeaf[]
      \nodepart{three}\ttfamily GetNode(SVOLink)\\ TryFindNodeIndex()\\ IsVoxelOccupied()
    };

    \node[clase] (node) at (0,-3.7) {
      {\itshape struct}\\[-1pt]\textbf{SVONode}
      \nodepart{two}\ttfamily MortonCode\\ FirstChild : SVOLink\\ Parent : SVOLink\\ Neighbors : NeighborSet
      \nodepart{three}\ttfamily HasChildren : bool
    };

    \node[clase] (leaf) at (4.7,-3.7) {
      {\itshape struct}\\[-1pt]\textbf{SVOLeaf}
      \nodepart{two}\ttfamily occupiedSubnodes : ulong
      \nodepart{three}\ttfamily IsOccupied(i)\\ IsEmpty / IsFull
    };

    \node[clase] (morton) at (-4.4,-7.6) {
      {\itshape struct}\\[-1pt]\textbf{MortonCode}
      \nodepart{two}\ttfamily code : uint
      \nodepart{three}\ttfamily ParentCode\\ ChildCode(i)\\ TryGetNeighborCode()
    };

    \node[clase] (nset) at (0,-7.6) {
      {\itshape struct}\\[-1pt]\textbf{NeighborSet}
      \nodepart{two}\ttfamily posX, negX, posY,\\ negY, posZ, negZ : SVOLink
      \nodepart{three}\ttfamily this[Direction]
    };

    \node[clase] (link) at (4.4,-7.6) {
      {\itshape struct}\\[-1pt]\textbf{SVOLink}
      \nodepart{two}\ttfamily link : uint
      \nodepart{three}\ttfamily IsNode(out layer)\\ IsVoxel(out sub)\\ Offset
    };

    % --- Relaciones (rombo en el contenedor, flecha discontinua para referencias) ---
    \draw[comp] (node.north) -- node[etiq, midway] {Layers} (svo.south);
    \draw[comp] (leaf.north) -- node[etiq, pos=0.62] {LeafNodes} (svo.east);
    \draw[comp] (morton.north) -- node[mult, pos=0.85, left] {1} (node.south west);
    \draw[comp] (nset.north) -- node[etiq, midway] {Neighbors} (node.south);
    \draw[comp] (link.west) -- node[mult, midway, above] {6} (nset.east);
    \draw[deref] (node.south east) -- node[etiq, pos=0.55] {FirstChild / Parent} (link.north);

  \end{tikzpicture}
  \caption[Estructuras de datos del espacio de nombres \texttt{Core} y sus relaciones.]
  {Estructuras de datos del espacio de nombres \texttt{Core} y sus relaciones.
  Los rombos denotan composición (propiedad de un valor) y las flechas discontinuas, referencias mediante enlaces.
  El \acrshort{svo} agrega las capas de \texttt{SVONode} y el vector paralelo de \texttt{SVOLeaf}; cada nodo referencia a su primer hijo y a su padre, y agrega su conjunto de seis vecinos, todos ellos codificados como \texttt{SVOLink} (el puntero compacto de la Sección~\ref{sec:svolink}).}
  \label{fig:core-clases}
\end{figure}

\section{Iteración 2: Construcción y rasterización del volumen}
\label{sec:iter-construccion}

Definida la estructura de datos, esta iteración aborda su construcción a partir de la geometría de una escena.
Este proceso, conocido como rasterización, consiste en consultar qué regiones del espacio están ocupadas por obstáculos para poblar el \acrshort{svo}.
Constituye una de las fases más costosas del sistema, por lo que su diseño está dominado por consideraciones de rendimiento.

La construcción se organiza como una canalización (\textit{pipeline}) de etapas bien diferenciadas, resumida en la Figura~\ref{fig:build-pipeline}: cada etapa consume el producto de datos de la anterior, y las dos etapas dominadas por consultas físicas se delegan en el \textit{Job System}.
Las subsecciones siguientes detallan estas etapas.

\begin{figure}[htp]
  \centering
  \begin{tikzpicture}[
      font=\small, node distance=5mm,
      etapa/.style={paso, text width=3.9cm, minimum height=1.0cm},
      tienda/.style={almacen, text width=3.9cm, minimum height=1.0cm},
      dato/.style={draw=udcgray!70!black, fill=udcgray!12, rounded corners=2pt,
                   align=left, font=\scriptsize, text width=3.0cm, inner sep=4pt},
      job/.style={draw=ficblue!60, fill=ficblue!8, rounded corners=2pt,
                  align=center, font=\scriptsize, text width=2.6cm, inner sep=4pt},
      conector/.style={densely dotted, draw=udcgray!75},
      conjob/.style={densely dotted, draw=ficblue!65}
    ]

    % --- Columna central: etapas de la canalizacion ---
    \node[tienda] (escena) {Escena\\[1pt]{\scriptsize(\textit{colliders})}};
    \node[etapa, below=of escena] (s1) {\textbf{1.~Rasterización L1}\\[1pt]{\scriptsize barrido grueso-a-fino}};
    \node[etapa, below=of s1] (s2) {\textbf{2.~Reserva de la capa 0}\\[1pt]{\scriptsize creación de nodos hoja}};
    \node[etapa, below=of s2] (s3) {\textbf{3.~Rasterización de hojas}\\[1pt]{\scriptsize descarte grueso + muestreo fino}};
    \node[etapa, below=of s3] (s4) {\textbf{4.~Capas superiores}\\[1pt]{\scriptsize de la raíz hacia las hojas}};
    \node[etapa, below=of s4] (s5) {\textbf{5.~Enlazado padre--hijos}};
    \node[etapa, below=of s5] (s6) {\textbf{6.~Enlazado de vecinos}};
    \node[tienda, below=of s6] (out) {SVO enlazado\\[1pt]{\scriptsize(\textit{NavContext})}};

    \foreach \a/\b in {escena/s1, s1/s2, s2/s3, s3/s4, s4/s5, s5/s6, s6/out}
      \draw[flecha] (\a) -- (\b);

    % --- Columna derecha: productos de datos ---
    \node[dato, right=9mm of s1] (d1) {códigos Morton de las celdas L1 ocupadas};
    \node[dato, right=9mm of s2] (d2) {\texttt{SVONode[]} de la capa 0};
    \node[dato, right=9mm of s3] (d3) {máscaras de 64\,bits (\texttt{SVOLeaf[]})};
    \node[dato, right=9mm of s4] (d4) {nodos internos por capa};
    \node[dato, right=9mm of s5] (d5) {enlaces \texttt{FirstChild} / \texttt{Parent}};
    \node[dato, right=9mm of s6] (d6) {\texttt{NeighborSet} (6 vecinos/nodo)};
    \foreach \s/\d in {s1/d1, s2/d2, s3/d3, s4/d4, s5/d5, s6/d6}
      \draw[conector] (\s.east) -- (\d.west);

    % --- Columna izquierda: etapas paralelizadas ---
    \node[job, left=9mm of s1] (j1) {\textit{Job System} + Burst\\[1pt]\texttt{OverlapBox} por lotes};
    \node[job, left=9mm of s3] (j3) {\textit{Job System} + Burst\\[1pt]descarte / fino / reducción};
    \draw[conjob] (j1.east) -- (s1.west);
    \draw[conjob] (j3.east) -- (s3.west);

  \end{tikzpicture}
  \caption[Canalización de construcción del \acrshort{svo}.]
  {Canalización de construcción del \acrshort{svo} durante el \gls{bake}.
  Cada etapa (centro) consume el producto de datos de la anterior (derecha);
  las etapas~1 y~3, dominadas por consultas físicas, se aceleran mediante el \textit{Job System} y Burst (izquierda) y se detallan en las Figuras~\ref{fig:coarse-to-fine} y~\ref{fig:empty-leaf}, mientras que el enlazado de la vecindad (etapa~6) se ilustra en la Figura~\ref{fig:neighbor-link}.}
  \label{fig:build-pipeline}
\end{figure}

\subsection{Muestreo del espacio mediante consultas físicas paralelizadas}

Tal y como se anticipó en los fundamentos, el sistema no lee directamente los datos de las mallas, sino que delega la detección de geometría en el motor de físicas a través de consultas \texttt{OverlapBox}, que comprueban si una caja determinada intersecta con algún \textit{collider}.
Dado el ingente número de consultas requeridas, estas se agrupan en lotes (\textit{batches}) mediante la estructura \texttt{OverlapBoxCommand} y se ejecutan en paralelo sobre varios hilos secundarios gracias al \textit{Job System} y al compilador Burst descritos en la Sección~\ref{sec:dots}.

La construcción se realiza con una estrategia \textbf{de lo grueso a lo fino} (\textit{coarse-to-fine}).
En lugar de muestrear de golpe la totalidad del espacio a máxima resolución, el barrido comienza en una capa cercana a la raíz y, en cada nivel, solo subdivide aquellas celdas que contienen geometría, descartando subárboles enteros vacíos (véase la Figura~\ref{fig:coarse-to-fine}).
Este descarte temprano es la primera de las dos grandes optimizaciones de la fase de construcción.

\begin{figure}[htp]
  \centering
  \begin{tikzpicture}[font=\small]
    % Izquierda: nivel grueso
    \begin{scope}
      \fill[obstaculo, rounded corners=2pt] (1.0,-1.0) rectangle (2.7,-2.7);
      \foreach \c in {0,1,2,3}\foreach \r in {0,1,2,3}
        \draw[celda, line width=0.7pt] (\c*0.8,-\r*0.8) rectangle ++(0.8,-0.8);
      \foreach \c/\r in {1/1,2/1,1/2,2/2}
        \draw[draw=udcpink, very thick] (\c*0.8,-\r*0.8) rectangle ++(0.8,-0.8);
      \node[font=\footnotesize, anchor=north, text width=3.4cm, align=center] at (1.6,-3.5)
        {Nivel grueso: se detectan las celdas con geometría};
    \end{scope}
    % Flecha
    \draw[flecha] (3.5,-1.6) -- ++(1.3,0) node[midway, above, font=\scriptsize] {subdivisión};
    % Derecha: nivel fino
    \begin{scope}[xshift=5.3cm]
      \fill[obstaculo, rounded corners=2pt] (1.0,-1.0) rectangle (2.7,-2.7);
      \foreach \c in {0,1,2,3}\foreach \r in {0,1,2,3}
        \draw[celda, line width=0.7pt] (\c*0.8,-\r*0.8) rectangle ++(0.8,-0.8);
      \foreach \c/\r in {1/1,2/1,1/2,2/2}
        \foreach \sc in {0,1}\foreach \sr in {0,1}
          \draw[celda] (\c*0.8+\sc*0.4,-\r*0.8-\sr*0.4) rectangle ++(0.4,-0.4);
      \foreach \c/\r in {1/1,2/1,1/2,2/2}
        \draw[draw=udcpink, very thick] (\c*0.8,-\r*0.8) rectangle ++(0.8,-0.8);
      \node[font=\footnotesize, anchor=north, text width=3.6cm, align=center] at (1.6,-3.5)
        {Nivel fino: solo esas celdas se subdividen};
    \end{scope}
  \end{tikzpicture}
  \caption[Barrido jerárquico de lo grueso a lo fino durante la rasterización.]
  {Barrido jerárquico de lo grueso a lo fino durante la rasterización.
  Las celdas vacías se descartan por completo; únicamente las que contienen geometría (resaltadas en rosa) se subdividen en el siguiente nivel de resolución.}
  \label{fig:coarse-to-fine}
\end{figure}

\subsection{Optimización: descarte de hojas vacías}
\label{sec:descarte-hojas}

La segunda optimización, y una de las más relevantes del sistema, se aplica durante la rasterización de las hojas.
Comprobar la ocupación de una hoja completa exige, en principio, lanzar 64 consultas \texttt{OverlapBox}, una por cada \gls{voxel}.
Sin embargo, este coste es casi siempre innecesario.

La razón es estructural: en un entorno de navegación tridimensional típico, el espacio navegable es \textbf{mayoritariamente aire libre}.
Los obstáculos (paredes, edificios, terreno) ocupan una fracción pequeña del volumen total, por lo que la inmensa mayoría de las hojas no contienen geometría alguna.
Esta es, precisamente, la observación que justifica el uso de una representación dispersa.

Para explotarla, la rasterización de hojas se divide en dos pasadas, ilustradas en la Figura~\ref{fig:empty-leaf} y esquematizadas en la Figura~\ref{fig:empty-leaf-pseudo}:

\begin{enumerate}
    \item \textbf{Pasada gruesa de descarte:} se lanza una \textbf{única} consulta \texttt{OverlapBox} que cubre la totalidad de la hoja. Esta caja se calcula como la caja envolvente exacta de los 64 vóxeles (expandidos por el radio del agente), de modo que es imposible que una hoja que no intersecte aquí contenga algún vóxel ocupado. Si la consulta no detecta nada, la hoja se deja vacía sin más.
    \item \textbf{Pasada fina:} únicamente las hojas supervivientes, es decir, aquellas en las que la pasada gruesa sí detectó geometría, pagan el coste de las 64 consultas individuales que determinan qué vóxeles concretos están ocupados.
\end{enumerate}

Dado que las supervivientes son una minoría, esta optimización reduce el número de consultas físicas en un factor cercano a 64 para la inmensa mayoría del volumen, que es vacío.
El resultado de la pasada fina se condensa en la máscara de 64 bits de cada hoja mediante un \textit{job} de reducción dedicado.

\begin{figure}[H]
\begin{lstlisting}
// --- Pasada gruesa: una consulta por hoja completa ---
lanzar  BuildCoarseLeafCommands   -> comandos_gruesos   (1 por hoja)
ejecutar OverlapBox (en lote)     -> resultados_gruesos
supervivientes = { i : resultados_gruesos[i] golpea algo }

si no hay supervivientes:
    todas las hojas quedan vacias  // caso muy comun

// --- Pasada fina: solo sobre las supervivientes ---
lanzar  BuildLeafCommands(supervivientes) -> comandos  (64 por superviviente)
ejecutar OverlapBox (en lote)             -> resultados
lanzar  ReduceLeafBitmasks                -> mascara[superviviente]

hojas[*] = Vacia
para cada j en supervivientes:
    hojas[j] = DesdeBits(mascara[j])
\end{lstlisting}
\caption{Rasterización de hojas en dos pasadas: descarte grueso y muestreo fino.}
\label{fig:empty-leaf-pseudo}
\end{figure}

\begin{figure}[htp]
  \centering
  \begin{tikzpicture}[font=\small]
    % Izquierda: pasada gruesa, 1 caja envolvente
    \begin{scope}
      \foreach \c in {0,1,2,3}\foreach \r in {0,1,2,3}
        \draw[celda] (\c*0.7,-\r*0.7) rectangle ++(0.7,-0.7);
      \draw[draw=ficblue, very thick, rounded corners=2pt] (-0.12,0.12) rectangle (2.92,-2.92);
      \node[font=\footnotesize, anchor=north, text width=3.4cm, align=center] at (1.4,-3.2)
        {Pasada gruesa:\\ \textbf{1} consulta envolvente};
    \end{scope}
    \draw[flecha] (3.4,-1.4) -- ++(1.2,0);
    % Derecha: pasada fina, 1 caja por voxel
    \begin{scope}[xshift=5cm]
      \foreach \c in {0,1,2,3}\foreach \r in {0,1,2,3} {
        \draw[celda] (\c*0.7,-\r*0.7) rectangle ++(0.7,-0.7);
        \draw[draw=ficblue, thick] (\c*0.7+0.07,-\r*0.7-0.07) rectangle ++(0.56,-0.56);
      }
      \node[font=\footnotesize, anchor=north, text width=3.7cm, align=center] at (1.4,-3.2)
        {Pasada fina:\\ \textbf{64} consultas individuales};
    \end{scope}
  \end{tikzpicture}
  \caption[Pasada gruesa de descarte frente a pasada fina de muestreo.]
  {Pasada gruesa de descarte frente a pasada fina de muestreo.
  La pasada gruesa resuelve la hoja entera con una única consulta envolvente; solo si esta detecta geometría se ejecuta la pasada fina, con una consulta por vóxel (64 en una hoja de $4\times4\times4$).}
  \label{fig:empty-leaf}
\end{figure}

\subsection{Enlazado de la vecindad}
\label{sec:enlazado-vecinos}

Una vez rasterizadas las hojas y construidas las capas superiores, resta una última etapa: enlazar cada nodo con sus seis vecinos (en las direcciones $\pm X$, $\pm Y$ y $\pm Z$), información imprescindible para que el buscador de rutas pueda navegar el grafo.

El enlazado se resuelve recorriendo las capas de la raíz hacia las hojas y aplicando una regla con \textbf{herencia del padre}.
Para cada nodo y dirección, se busca primero si existe un vecino del mismo tamaño (mismo nivel) mediante el código de Morton correspondiente.
Si no existe, en lugar de buscar exhaustivamente vecinos de distinto tamaño, el nodo simplemente \textbf{hereda el enlace del vecino de su padre} (un nodo de mayor tamaño).
Esta solución, ilustrada en la Figura~\ref{fig:neighbor-link}, resuelve de forma elegante y eficiente la transición entre regiones de distinta resolución, que es uno de los aspectos más delicados de la navegación sobre estructuras jerárquicas.

\begin{figure}[htp]
  \centering
  \begin{tikzpicture}[font=\small]
    % Vecino grande (nivel del padre)
    \draw[celda, very thick, fill=ficblue!12] (0,0) rectangle (2.6,-2.6);
    \node[font=\footnotesize, align=center] at (1.3,-1.85) {Vecino de\\mayor tamaño};
    % Nodo a enlazar, subdividido
    \begin{scope}[xshift=2.6cm]
      \draw[celda, very thick] (0,0) rectangle (2.6,-2.6);
      \draw[celda] (1.3,0) -- (1.3,-2.6);
      \draw[celda] (0,-1.3) -- (2.6,-1.3);
      \foreach \sc in {0,1}\foreach \sr in {0,1}
        \draw[celda] (\sc*0.65,-\sr*0.65) rectangle ++(0.65,-0.65);
      \fill[udcpink!65] (0,0) rectangle (0.65,-0.65);
      \draw[celda] (0,0) rectangle (0.65,-0.65);
    \end{scope}
    \node[font=\scriptsize, anchor=west] at (3.4,-0.32) {nodo pequeño};
    % Enlace heredado del padre
    \draw[flecha, draw=udcpink, very thick] (2.6,-0.32) to[out=180, in=0] (1.6,-0.78);
    \node[font=\scriptsize, text=udcpink!80!black, anchor=north] at (1.5,-1.0) {vecino heredado};
  \end{tikzpicture}
  \caption[Enlazado de la vecindad mediante herencia del nodo padre.]
  {Enlazado de la vecindad mediante herencia del nodo padre.
  Cuando un nodo pequeño carece de un vecino de su mismo tamaño en una dirección, hereda como vecino el nodo de mayor tamaño adyacente a su padre, evitando una búsqueda exhaustiva entre celdas de distinto tamaño.}
  \label{fig:neighbor-link}
\end{figure}

Al término de esta iteración el sistema es ya capaz de transformar una escena tridimensional en un \acrshort{svo} completamente enlazado y navegable.

\section{Iteración 3: Búsqueda de rutas tridimensional}
\label{sec:iter-pathfinding}

Esta iteración implementa el corazón del sistema y completa el producto mínimo viable: la búsqueda de una trayectoria libre de colisiones entre dos puntos.
Conforme a lo expuesto en los fundamentos, se adapta el algoritmo \gls{astar}~\cite{astar} para operar sobre el \acrshort{svo}.

\subsection{Búsqueda sobre un grafo de resolución heterogénea}

La particularidad esencial de este buscador, frente a una implementación clásica de \gls{astar} sobre una rejilla uniforme, es que el grafo subyacente posee una \textbf{resolución heterogénea}: un mismo salto puede atravesar un nodo enorme de espacio abierto o un diminuto \gls{voxel} junto a un obstáculo.
El algoritmo gestiona esta heterogeneidad combinando tres comportamientos durante la expansión de vecinos, ilustrados en la Figura~\ref{fig:astar-multinivel}:

\begin{itemize}
    \item Si un vecino es un nodo grande de \textbf{espacio libre}, se trata como un único salto, lo que permite cruzar grandes extensiones vacías con muy pocas expansiones.
    \item Si un vecino es un nodo que \textbf{contiene geometría}, no puede tratarse como un salto barato y seguro: se expande en sus ocho hijos para examinarlos a mayor resolución.
    \item Si un vecino es una \textbf{hoja parcialmente ocupada}, se expande directamente en sus \glspl{voxel} libres, descartando los ocupados.
\end{itemize}

\begin{figure}[htp]
  \centering
  \begin{tikzpicture}[font=\small]
    % Marco exterior
    \draw[celda, line width=0.7pt] (0,0) rectangle (6.4,-3.6);
    % Zona abierta (izquierda): pocas celdas grandes
    \draw[celda, line width=0.7pt] (1.8,0) -- (1.8,-3.6);
    \draw[celda, line width=0.7pt] (3.6,0) -- (3.6,-3.6);
    \draw[celda, line width=0.7pt] (0,-1.8) -- (3.6,-1.8);
    % Zona fina (derecha, cerca del obstaculo)
    \foreach \c in {0,...,3}\foreach \r in {0,...,7}
      \draw[celda] (3.6+\c*0.7,-\r*0.45) rectangle ++(0.7,-0.45);
    % Obstaculo
    \fill[obstaculo, rounded corners=2pt] (4.2,-1.0) rectangle (5.6,-2.7);
    \node[font=\footnotesize, text=white] at (4.9,-1.85) {obstáculo};
    % Ruta A*
    \node[nodo, fill=ficblue!40] (s) at (0.7,-0.9) {};
    \node[font=\scriptsize, anchor=south] at (0.7,-0.62) {inicio};
    \coordinate (a) at (2.7,-0.9);
    \coordinate (b) at (3.9,-0.7);
    \coordinate (c) at (5.0,-0.55);
    \coordinate (d) at (5.95,-1.2);
    \coordinate (e) at (5.95,-2.6);
    \coordinate (f) at (5.2,-3.15);
    \node[nodo, fill=udcpink!75] (g) at (4.4,-3.15) {};
    \node[font=\scriptsize, anchor=north] at (4.4,-3.35) {meta};
    \draw[flecha, draw=udcpink, very thick] (s) -- (a) -- (b) -- (c) -- (d) -- (e) -- (f) -- (g);
    \foreach \p in {a,b,c,d,e,f} \fill[udcpink] (\p) circle (1.7pt);
    \node[font=\scriptsize, text=udcpink!80!black, anchor=south] at (1.7,-0.85) {salto grande};
    \node[font=\scriptsize, text=udcpink!80!black, anchor=west] at (6.0,-1.95) {pasos finos};
  \end{tikzpicture}
  \caption[Expansión de A* sobre un grafo de resolución heterogénea.]
  {Expansión de A* sobre un grafo de resolución heterogénea.
  En el espacio abierto la búsqueda avanza con saltos grandes entre nodos de gran tamaño; al aproximarse al obstáculo, se ramifica en nodos y vóxeles cada vez más pequeños para bordearlo con precisión.}
  \label{fig:astar-multinivel}
\end{figure}

\begin{figure}[H]
\begin{lstlisting}
empujar(inicio, f = h(inicio))
mientras lista_abierta no vacia:
    n = extraer_minimo()
    si cerrado[n]: continuar
    marcar cerrado[n]
    si n == meta: devolver reconstruir(n)

    para cada m en vecinos(n):
        si m es nodo grande con geometria: expandir sus hijos; continuar
        si m es hoja parcialmente ocupada: expandir sus voxeles libres; continuar
        si bloqueado(m): continuar

        g_nuevo = g[n] + coste(n, m)
        si g_nuevo < g[m]:
            g[m] = g_nuevo;  procedencia[m] = n
            empujar(m, g_nuevo + h(m))
\end{lstlisting}
\caption{Esquema del bucle principal de la búsqueda \gls{astar} adaptada.}
\label{fig:astar-pseudo}
\end{figure}

\subsection{Función de coste y heurística}

El sistema ofrece dos modos de medir el coste de la trayectoria.
El modo basado en \textbf{distancia euclídea} computa el coste real entre los centros de los nodos, produciendo rutas geométricamente más cortas.
El modo basado en \textbf{conteo de nodos}, en cambio, asigna un coste unitario a cada salto con independencia del tamaño del nodo; al penalizar por igual cualquier salto, este modo sesga la búsqueda hacia los nodos grandes y, por tanto, hacia el espacio abierto, alejando al agente de los obstáculos.

La heurística empleada es la distancia en línea recta hasta la meta, ponderada por un peso configurable $W$ (heurística voraz).
Para el modo de conteo de nodos, dado que el coste $g$ no se expresa en unidades de distancia, la heurística se reescala dividiéndola por el tamaño del nodo más grande, garantizando así que ambos términos de la función $f(n) = g(n) + W \cdot h(n)$ sean comparables.

\subsection{Reconstrucción y robustez}

Cuando la búsqueda alcanza la meta, la trayectoria se reconstruye recorriendo hacia atrás los enlaces de procedencia almacenados para cada nodo.
El buscador contempla además los casos límite, como volúmenes completamente vacíos (en los que la ruta es trivial), extremos no válidos o situados sobre geometría, y un presupuesto máximo de nodos que acota el coste de las búsquedas sin solución.

Al completar esta iteración se dispone del producto mínimo viable comprometido en la planificación: un sistema capaz de construir el volumen y calcular una ruta sobre él.
No obstante, tal y como se documentó en el seguimiento de la planificación (Sección~\ref{sec:seguimiento}), garantizar la robustez geométrica de estas rutas resultaría más costoso de lo previsto, cuestión que se aborda en la Iteración~6.

\section{Iteración 4: Postprocesado de trayectorias}
\label{sec:iter-postprocesado}

Las rutas producidas por la búsqueda sobre el volumen discretizado resultan artificiales, al estar compuestas por multitud de segmentos alineados con la rejilla.
Esta iteración implementa el postprocesado en dos fases descrito en los fundamentos, con el fin de dotar al agente de un movimiento orgánico.

\subsection{Atajo voraz mediante línea de visión}

La primera fase aplica un algoritmo de \textit{string pulling} en su variante voraz, que elimina los nodos intermedios redundantes.
Partiendo del primer punto, el algoritmo intenta conectarlo con el punto más lejano posible siempre que exista \textbf{línea de visión directa} entre ambos, descartando los nodos intermedios.

La comprobación de la línea de visión es un punto crítico de precisión y rendimiento.
Se resuelve mediante el algoritmo \acrshort{dda}~\cite{amanatides1987fast}, que recorre de forma incremental los \glspl{voxel} atravesados por el segmento, comprobando si alguno está ocupado.
El recorrido avanza, en cada paso, por el eje cuyo siguiente plano de \gls{voxel} se encuentra más próximo, evitando así tanto omitir \glspl{voxel} como evaluarlos por duplicado.
La Figura~\ref{fig:dda-pseudo} resume el procedimiento.

\begin{figure}[H]
\begin{lstlisting}
celda   = floor(inicio)              // voxel inicial
paso    = signo(direccion)           // +1 o -1 por eje
tDelta  = 1 / |direccion|            // avance parametrico por eje
tCruce  = distancia al primer plano de voxel por eje

mientras dentro de la rejilla:
    si ocupado(celda): devolver false      // obstaculo: sin linea de vision

    eje = aquel con menor tCruce            // siguiente plano mas cercano
    si tCruce[eje] >= longitud_segmento: terminar
    celda[eje]  += paso[eje]
    tCruce[eje] += tDelta[eje]

devolver true
\end{lstlisting}
\caption{Recorrido \acrshort{dda} para la comprobación de línea de visión.}
\label{fig:dda-pseudo}
\end{figure}

\subsection{Suavizado mediante \textit{splines} de Catmull-Rom}

La trayectoria depurada se suaviza interpolando sus puntos con un \textit{spline} de Catmull-Rom, evaluado a una resolución fija entre cada par de puntos de control.
Como se argumentó en los fundamentos (véase la Figura~\ref{fig:catmull_rom}), la elección de esta curva se debe a que \textbf{pasa de forma exacta por todos los puntos de control}.
Esta propiedad es esencial para la seguridad de la trayectoria: dado que la fase anterior garantiza que dichos puntos se hallan en espacio libre, obligar a la curva a cruzar exactamente por ellos evita que el suavizado introduzca colisiones accidentales con la geometría.
La evaluación de la curva se muestra en la Figura~\ref{fig:catmull-eval}.

\begin{figure}[H]
\begin{lstlisting}[language={[Sharp]C}]
static Vector3 EvaluateCatmullRom(Vector3 p0, Vector3 p1, Vector3 p2, Vector3 p3, float t) =>
    0.5f * (
        2f * p1
        + (-p0 + p2) * t
        + (2f * p0 - 5f * p1 + 4f * p2 - p3) * (t * t)
        + (-p0 + 3f * p1 - 3f * p2 + p3) * (t * t * t)
    );
\end{lstlisting}
\caption{Evaluación de un segmento del \textit{spline} de Catmull-Rom.}
\label{fig:catmull-eval}
\end{figure}

\section{Iteración 5: Persistencia e integridad de la escena}
\label{sec:iter-persistencia}

Reconstruir el \acrshort{svo} cada vez que se carga la escena sería prohibitivo en tiempo de ejecución.
Esta iteración implementa el proceso de \gls{bake} descrito en los fundamentos: la construcción anticipada del volumen desde el editor y su almacenamiento en disco para una carga posterior inmediata.

\subsection{Serialización en un único bloque binario}

El \acrshort{svo} se serializa en un \textbf{único bloque binario} (\textit{blob}) que se almacena dentro de un \textit{asset} del motor.
Cada hoja se vuelca como su entero de 64 bits, y cada nodo como nueve enteros de 32 bits (su código de Morton, los enlaces a su primer hijo y a su padre, y sus seis vecinos).
Esta representación plana y sin referencias es a la vez compacta y extremadamente rápida de cargar, ya que la reconstrucción se reduce a una lectura secuencial del flujo de bytes, sin necesidad de recalcular la vecindad ni la jerarquía.

\subsection{Integridad de la escena: reinterpretación de los bits de un \textit{float}}
\label{sec:hash-float}

Un riesgo evidente del \gls{bake} es que la escena se modifique tras él, dejando los datos almacenados obsoletos.
Para detectarlo, el sistema calcula un \textit{hash} de la geometría relevante en el momento del \gls{bake} y lo compara con el de la escena actual al cargar.
Conforme se justificó en los fundamentos, se emplea el algoritmo \gls{fnv1a}~\cite{fnv_hash} en lugar del \texttt{GetHashCode} por defecto, ya que este último no es determinista entre ejecuciones~\cite{csharp_gethashcode}.

El algoritmo \gls{fnv1a} opera sobre los \textbf{bytes} de los datos, lo que plantea un reto al alimentarlo con valores en coma flotante (las posiciones, rotaciones y escalas de los obstáculos): es necesario obtener el patrón de bits \textbf{exacto} de cada \texttt{float}, y no una conversión numérica que truncaría o redondearía su valor.
Sin embargo, C\# no permite reinterpretar directamente los bits de un tipo como los de otro.

La solución implementada recurre a una \textbf{unión} al estilo del lenguaje C, ilustrada en la Figura~\ref{fig:float-union}.
Mediante el atributo \texttt{[StructLayout(LayoutKind.Explicit)]} se define una estructura con dos campos, un \texttt{float} y un \texttt{int}, forzados a compartir la misma posición de memoria con \texttt{[FieldOffset(0)]}.
Al escribir el valor en el campo \texttt{float} y leer acto seguido el campo \texttt{int}, se obtiene el entero cuya representación binaria coincide bit a bit con la del \texttt{float} original.
Esta técnica es eficiente (no realiza reservas de memoria ni saltos condicionales) y garantiza el determinismo necesario para que el \textit{hash} sea fiable entre distintas sesiones.

\begin{figure}[H]
\begin{lstlisting}[language={[Sharp]C}]
[StructLayout(LayoutKind.Explicit)]
struct FloatIntUnion
{
    [FieldOffset(0)] public float Float;
    [FieldOffset(0)] public int   Int;
}

// Devuelve el int cuya representacion de bits coincide con la del float.
public static int FloatBitsAsInt(float value) =>
    new FloatIntUnion { Float = value }.Int;
\end{lstlisting}
\caption{Reinterpretación de los bits de un \texttt{float} como un \texttt{int} mediante una unión.}
\label{fig:float-union}
\end{figure}

Para que el \textit{hash} no dependa del orden arbitrario en que el motor devuelve los \textit{colliders}, estos se ordenan previamente por una huella estable (posición, rotación y escala) antes de alimentar el algoritmo.

\section{Iteración 6: Robustez geométrica y radio del agente}
\label{sec:iter-robustez}

Esta iteración aborda directamente la desviación más significativa de la planificación, anticipada en la Sección~\ref{sec:seguimiento}: la consolidación de la robustez geométrica de las rutas.
Las primeras versiones del buscador producían, en ocasiones, trayectorias que rozaban o atravesaban la geometría del entorno (un fenómeno conocido como \textit{clipping}), especialmente en las inmediaciones de las aristas y esquinas de los obstáculos.

\subsection{Consideración del radio del agente}

La causa principal residía en que el sistema trataba al agente como un punto adimensional.
La solución consiste en \textbf{inflar} los obstáculos por el radio del agente durante la rasterización: al expandir cada caja de consulta \texttt{OverlapBox} por dicho radio, todo \gls{voxel} situado a una distancia de la geometría inferior al tamaño del agente se marca como ocupado.
De este modo, cualquier ruta que el buscador considere libre garantiza, por construcción, el espacio de holgura necesario para que el agente no colisione.
La Figura~\ref{fig:clipping} ilustra el problema y su corrección.

\begin{figure}[htp]
  \centering
  \begin{tikzpicture}[font=\small]
    % ===== Izquierda: sin tratar el radio (colision) =====
    \begin{scope}
      \fill[obstaculo, rounded corners=1pt] (1.6,-1.2) rectangle (3.0,-2.6);
      \node[font=\footnotesize, text=white] at (2.3,-1.9) {obstáculo};
      \draw[flecha, draw=udcpink, very thick] (0,-0.2) -- (3.6,-2.05);
      \fill[ficblue!25, opacity=0.75] (1.85,-1.05) circle (0.34);
      \draw[agente] (1.85,-1.05) circle (0.34);
      \node[font=\scriptsize, text=udcpink!80!black, anchor=south] at (1.1,-0.32) {colisión};
      \draw[densely dotted, thin, draw=udcpink!80!black] (1.1,-0.42) -- (1.72,-1.0);
      \node[font=\footnotesize, anchor=north, text width=3.9cm, align=center] at (1.8,-3.05)
        {Sin tratar el radio: la trayectoria recorta la esquina};
    \end{scope}
    % ===== Derecha: obstaculo inflado (holgura) =====
    \begin{scope}[xshift=5.4cm]
      \draw[dashed, draw=udcpink, thick, rounded corners=4pt] (1.3,-0.9) rectangle (3.3,-2.9);
      \fill[obstaculo, rounded corners=1pt] (1.6,-1.2) rectangle (3.0,-2.6);
      \node[font=\footnotesize, text=white] at (2.3,-1.9) {obstáculo};
      \draw[flecha, draw=ficblue, very thick] (0,-0.2) -- (1.15,-0.5) -- (3.45,-0.5) -- (3.8,-2.0);
      \draw[<->, thin, draw=udcpink!80!black] (2.3,-0.5) -- (2.3,-0.9);
      \node[font=\scriptsize, text=udcpink!80!black, anchor=west] at (2.4,-0.7) {holgura $=$ radio};
      \node[font=\footnotesize, anchor=north, text width=4.0cm, align=center] at (1.9,-3.25)
        {Obstáculo inflado por el radio: trayectoria con holgura};
    \end{scope}
  \end{tikzpicture}
  \caption[Corrección del \textit{clipping} mediante el inflado de obstáculos por el radio del agente.]
  {Corrección del \textit{clipping} mediante el inflado de obstáculos por el radio del agente.
  Si no se considera el tamaño del agente, una trayectoria que recorta una esquina provoca colisión (izquierda).
  Al inflar el obstáculo por el radio del agente (margen discontinuo), la trayectoria conserva siempre la holgura necesaria (derecha).}
  \label{fig:clipping}
\end{figure}

\subsection{Validación intensiva mediante pruebas}

Diagnosticar de forma fiable este tipo de errores geométricos, intermitentes y dependientes de la configuración concreta de la escena, exigió el desarrollo de una batería de pruebas considerablemente más extensa de la prevista.
A las pruebas en modo de edición se sumaron pruebas en modo de ejecución (\textit{PlayMode}), capaces de instanciar escenas completas, construir el volumen y validar la navegabilidad de las trayectorias resultantes contra el motor de físicas real.
Tanto estas pruebas en modo de ejecución como las pruebas en modo de edición existentes se ejecutan de forma automática en el sistema de \acrshort{ci} ante cada \textit{pull request}, de modo que cualquier regresión queda detectada antes de integrar los cambios en la rama principal.
Este esfuerzo de validación, aunque desplazó ligeramente el calendario, resultó determinante para consolidar un núcleo fiable sobre el que sustentar el resto de funcionalidades.

\section{Iteración 7: Integración en Unity y telemetría}
\label{sec:iter-integracion}

Con un núcleo robusto, esta iteración se centra en la usabilidad, materializando el sistema en herramientas accesibles tanto para perfiles técnicos como para diseñadores y artistas.

\subsection{Componentes y modos de construcción}

El sistema se expone principalmente a través del componente \texttt{NavVolumeSpace}, que encapsula un volumen de navegación y ofrece tres modos de construcción para adaptarse a distintos flujos de trabajo: \textbf{precalculado} (\textit{Baked}), que carga los datos serializados; \textbf{construido al inicio} (\textit{BuildOnAwake}); y \textbf{manual}.
A este se añaden los componentes que representan a las entidades de la escena, como los agentes y los obstáculos dinámicos.

\subsection{Visualización mediante Gizmos}

Para facilitar la comprensión y la depuración del sistema, se desarrolló una visualización del \acrshort{svo} en la vista de escena mediante Gizmos.
Esta representa la jerarquía del árbol con un código de colores que distingue los nodos vacíos, los que contienen geometría, las hojas parciales y los \glspl{voxel} ocupados (véase la Figura~\ref{fig:gizmos}).
Dado que un volumen puede contener cientos de miles de nodos, el dibujado incorpora dos salvaguardas de rendimiento: un \textbf{presupuesto máximo} de primitivas y un \textbf{descarte por distancia} a la cámara, que evita representar el detalle fino de regiones lejanas.

\pendientefig{Captura de la vista de escena de Unity mostrando la visualización del \acrshort{svo} mediante Gizmos sobre un escenario con obstáculos: nodos de distintos tamaños coloreados según su estado (vacíos, con geometría, hojas parciales) y los vóxeles ocupados resaltados.}{Visualización del \acrshort{svo} en la vista de escena mediante Gizmos.}{fig:gizmos}

\subsection{Telemetría: el cálculo de la memoria estimada}
\label{sec:memoria-estimada}

El inspector incorpora un panel de estadísticas que ayuda al usuario a dimensionar y ajustar el volumen (Figura~\ref{fig:stats-panel}).
Entre las métricas más relevantes figuran el ahorro de \glspl{voxel} frente a una representación densa y una \textbf{estimación de la memoria} ocupada por la estructura.

El cálculo de esta memoria estimada se fundamenta en los principios de disposición de objetos en la memoria administrada de C\#.
En una plataforma de 64 bits, todo objeto reservado en la \gls{heap} acarrea una sobrecarga: una cabecera de objeto (que incluye el puntero a su tabla de métodos y el bloque de sincronización) y, en el caso de los vectores, espacio adicional para almacenar su longitud.
A partir de estas constantes aproximadas, la estimación recorre la estructura sumando, para cada vector de nodos y de hojas, su cabecera más el tamaño real de cada elemento (obtenido con \texttt{Unsafe.SizeOf}) multiplicado por su número de elementos, tal y como se muestra en la Figura~\ref{fig:memoria-pseudo}.

\begin{figure}[H]
\begin{lstlisting}[language={[Sharp]C}]
const int _OBJECT_HEADER_BYTES = 24; // cabecera de objeto
const int _ARRAY_HEADER_BYTES  = 24; // cabecera de objeto + longitud
const int _REFERENCE_BYTES     = 8;  // tamano de una referencia

int total = _OBJECT_HEADER_BYTES;

// Vector de hojas: cada hoja es un ulong (8 bytes).
total += _ARRAY_HEADER_BYTES + Unsafe.SizeOf<SVOLeaf>() * leaves.Length;

// Vector externo de capas (un vector dentado de referencias).
total += _ARRAY_HEADER_BYTES + _REFERENCE_BYTES * layers.Length;

// Cada capa: su cabecera mas el tamano real de cada nodo.
foreach (var layer in layers)
    total += _ARRAY_HEADER_BYTES + Unsafe.SizeOf<SVONode>() * layer.Length;
\end{lstlisting}
\caption{Estimación de la memoria ocupada por el \acrshort{svo}.}
\label{fig:memoria-pseudo}
\end{figure}

Conviene subrayar que se trata de una \textbf{estimación}, no de una medición exacta: su propósito es ofrecer al usuario una magnitud orientativa y, sobre todo, comparable entre configuraciones.
Por su parte, el ahorro de espacio se cuantifica comparando el número de \glspl{voxel} efectivamente reservados frente al de un árbol completamente denso de la misma profundidad, que crece como $8^{(L-1)}$ hojas para $L$ capas.
Esta cifra evidencia de forma tangible el beneficio de la representación dispersa, que constituye el fundamento de todo el sistema.

\pendientefig{Captura del panel de estadísticas del inspector: tamaño de vóxel, vóxeles reservados frente al total teórico con su porcentaje de ahorro, memoria estimada, y los gráficos de nodos por capa y de ahorro disperso por capa (mapa de calor).}{Panel de estadísticas del volumen en el inspector del editor.}{fig:stats-panel}

\section{Iteración 8: API pública, optimización del rendimiento}
\label{sec:iter-api-opt}

Esta iteración consolida la interfaz pública del sistema y acomete la fase de evaluación y optimización del rendimiento prevista en la planificación.

\subsection{API de navegación y consultas espaciales}

Sobre el componente \texttt{NavVolumeSpace} se expone una \acrshort{api} que cubre los casos de uso habituales de la navegación.
Más allá del cálculo de rutas, se ofrecen consultas espaciales como la comprobación de si un punto es navegable, el ajuste de un punto al \gls{voxel} navegable más cercano (\textit{snapping}) o la obtención de puntos navegables aleatorios, de utilidad para poblar escenas o dirigir el deambular de los agentes.

El cálculo de rutas se ofrece en dos variantes.
La variante \textbf{síncrona} reutiliza una única instancia del buscador, mientras que la variante \textbf{asíncrona} ejecuta la búsqueda en un hilo secundario, evitando bloquear el hilo principal en escenarios con muchas peticiones.
Esta última emplea una reserva (\textit{pool}) de buscadores para no incurrir en reservas de memoria repetidas, y admite la cancelación de la petición mediante un \textit{token}.
La Figura~\ref{fig:seq-path} detalla la secuencia de una petición asíncrona, desde la llamada del agente hasta la entrega de la trayectoria suavizada.

\begin{figure}[htp]
  \centering
  \resizebox{\textwidth}{!}{%
  \begin{tikzpicture}[
      font=\small,
      head/.style={draw=ficblue!70!black, fill=ficblue!12, thick, rounded corners=2pt,
                   align=center, font=\footnotesize, minimum height=8mm, inner xsep=5pt},
      life/.style={dashed, draw=black!45},
      msg/.style={-{Stealth[length=2.4mm]}, thick},
      ret/.style={-{Stealth[length=2.4mm]}, dashed, draw=black!60},
      lbl/.style={font=\scriptsize, inner sep=1.5pt, fill=white},
      act/.style={fill=ficblue!22, draw=ficblue!70!black, line width=0.3pt}
    ]

    % Coordenadas X de las lineas de vida
    \def\xA{0}     % Agente
    \def\xN{3.6}   % NavVolumeSpace
    \def\xP{6.7}   % Pool
    \def\xS{9.1}   % SVOPathfinder
    \def\xM{11.9}  % PathSmoother
    \def\ybot{-9.6}

    % Cabeceras
    \node[head] (hA) at (\xA,0) {Agente};
    \node[head] (hN) at (\xN,0) {NavVolumeSpace};
    \node[head] (hP) at (\xP,0) {Pool de\\buscadores};
    \node[head] (hS) at (\xS,0) {SVOPathfinder\\(A*)};
    \node[head] (hM) at (\xM,0) {PathSmoother};

    % Lineas de vida
    \foreach \x in {\xA,\xN,\xP,\xS,\xM}
      \draw[life] (\x,-0.45) -- (\x,\ybot);

    % Marco del trabajo en hilo secundario (capa normal, antes de barras y mensajes)
    \draw[draw=udcgray!70, thick, rounded corners=2pt] (\xN-1.05,-2.2) rectangle (\xM+1.05,-8.5);

    % Barras de activacion
    \filldraw[act] (\xN-0.09,-1.0)  rectangle (\xN+0.09,-9.1);
    \filldraw[act] (\xP-0.09,-2.8)  rectangle (\xP+0.09,-3.3);
    \filldraw[act] (\xS-0.09,-3.9)  rectangle (\xS+0.09,-5.3);
    \filldraw[act] (\xM-0.09,-5.9)  rectangle (\xM+0.09,-6.4);
    \filldraw[act] (\xM-0.09,-7.0)  rectangle (\xM+0.09,-7.5);
    \filldraw[act] (\xP-0.09,-8.1)  rectangle (\xP+0.09,-8.3);

    % Mensajes
    \draw[msg] (\xA,-1.0) -- (\xN,-1.0) node[lbl, midway, above] {FindPathAsync(request, token)};
    \draw[ret] (\xN,-1.7) -- (\xA,-1.7) node[lbl, midway, above] {Task\textless PathResult\textgreater};

    \draw[msg] (\xN,-2.8) -- (\xP,-2.8) node[lbl, midway, above] {TryTake()};
    \draw[ret] (\xP,-3.3) -- (\xN,-3.3) node[lbl, midway, above] {buscador (o nuevo)};

    \draw[msg] (\xN,-3.9) -- (\xS,-3.9) node[lbl, midway, above] {FindPath(ctx, request)};
    % Bucle interno de A*
    \draw[msg] (\xS+0.09,-4.5) -- ++(0.7,0) -- ++(0,-0.4) -- (\xS+0.09,-4.9);
    \node[lbl, anchor=west] at (\xS+0.85,-4.7) {bucle A*};
    \draw[ret] (\xS,-5.3) -- (\xN,-5.3) node[lbl, midway, above] {ruta en bruto};

    \draw[msg] (\xN,-5.9) -- (\xM,-5.9) node[lbl, midway, above] {GreedyShortcut(raw)};
    \draw[ret] (\xM,-6.4) -- (\xN,-6.4) node[lbl, midway, above] {atajo (LOS / DDA)};

    \draw[msg] (\xN,-7.0) -- (\xM,-7.0) node[lbl, midway, above] {CatmullRomSpline(atajo)};
    \draw[ret] (\xM,-7.5) -- (\xN,-7.5) node[lbl, midway, above] {ruta suavizada};

    \draw[msg] (\xN,-8.1) -- (\xP,-8.1) node[lbl, midway, above] {Add(buscador)};

    \draw[ret] (\xN,-9.1) -- (\xA,-9.1) node[lbl, midway, above] {PathResult};

    % Etiqueta del fragmento, dibujada al final (opaca) para no quedar bajo la barra de NavVolumeSpace
    \node[anchor=north west, fill=white, draw=udcgray!70, rounded corners=1pt,
          inner xsep=4pt, inner ysep=2pt, font=\scriptsize\itshape, text=udcgray!85!black]
      at (\xN-1.05,-2.2) {hilo secundario (\texttt{Task.Run})};

  \end{tikzpicture}%
  }
  \caption[Secuencia de una petición de ruta asíncrona.]
  {Secuencia de una petición de ruta asíncrona.
  Tras devolver de inmediato una \texttt{Task}, el cálculo prosigue en un hilo secundario: se toma un buscador de la reserva, se ejecuta la búsqueda \gls{astar}, se postprocesa la ruta (atajo por línea de visión y suavizado por \textit{spline}) y se devuelve el buscador a la reserva antes de entregar el resultado.
  La variante síncrona reutiliza un único buscador y prescinde del hilo secundario.}
  \label{fig:seq-path}
\end{figure}

\subsection{Optimizaciones de rendimiento}

La fase de optimización, guiada por el perfilado del sistema, introdujo numerosas mejoras de bajo nivel coherentes con los fundamentos del \acrshort{dod}.
Entre las más significativas cabe destacar:

\begin{itemize}
    \item \textbf{Reutilización de \textit{buffers}:} el buscador reutiliza las colecciones de su estado interno entre peticiones, en lugar de reservarlas de nuevo, reduciendo la presión sobre el recolector de basura.
    \item \textbf{Comprobación de cancelación amortizada:} en lugar de consultar el \textit{token} de cancelación en cada expansión (una operación relativamente costosa), solo se comprueba periódicamente mediante una máscara de bits.
    \item \textbf{Acceso por referencia de solo lectura:} la consulta de nodos devuelve una referencia (\texttt{ref readonly}) en lugar de una copia de la estructura, evitando duplicar datos en las rutas más calientes.
    \item \textbf{Comparación sin \textit{boxing}:} como se adelantó en la Sección~\ref{sec:svolink}, los enlaces implementan \texttt{IEquatable} para su uso eficiente como clave en las colecciones \textit{hash} del buscador.
\end{itemize}

Asimismo, se optimizó la propia construcción del volumen, evitando reordenaciones innecesarias de las capas cuando no faltan nodos hermanos, lo que aceleró notablemente el \gls{bake} en escenas densas.

\section{Iteración 9: Evasión local entre agentes}
\label{sec:iter-evasion}

La última iteración incorpora un incremento que, tal y como se documentó en el seguimiento de la planificación (Sección~\ref{sec:seguimiento}), no figuraba como objetivo prioritario inicial, pero que la holgura recuperada permitió abordar: la evasión local entre agentes.
La búsqueda de rutas descrita hasta ahora calcula trayectorias sobre la geometría estática, pero no evita que varios agentes en movimiento colisionen entre sí.

\subsection{El algoritmo ORCA}

Para resolverlo se implementó el algoritmo \acrshort{orca}~\cite{van2011reciprocal}, introducido en los fundamentos, que razona sobre el espacio de velocidades para calcular, en cada instante, una velocidad segura y próxima a la deseada.
La implementación es un port de las rutinas de programación lineal en tres dimensiones de la biblioteca de referencia RVO2-3D~\cite{rvo2_3d}.
El solucionador busca la velocidad óptima que satisface todas las restricciones (semiplanos de velocidad) y, cuando el problema resulta infactible (demasiadas restricciones simultáneas), recurre a una relajación que minimiza la penetración en las restricciones (véase la Figura~\ref{fig:orca}).
Dado que este solucionador se ejecuta para cada agente y en cada fotograma, sus estructuras temporales se reservan en el \gls{stack} mediante \texttt{stackalloc} y se manipulan a través de \texttt{Span\textless T\textgreater}, lo que evita reservas en el \gls{heap} y la consiguiente presión sobre el recolector de basura en una de las rutas más críticas del sistema.

Una particularidad de la integración es que el sistema no solo considera la evasión recíproca entre agentes, sino que también deriva restricciones a partir de los \glspl{voxel} ocupados del \acrshort{svo} cercano, de modo que la evasión local respeta igualmente la geometría estática del entorno.

\begin{figure}[htp]
  \centering
  \begin{tikzpicture}[font=\small]
    % Region de velocidades prohibidas (semiplano sombreado)
    \fill[udcpink!16] (-2.6,0.1) -- (2.4,2.6) -- (2.8,2.8) -- (-2.6,2.8) -- cycle;
    % Ejes del espacio de velocidades
    \draw[->, thick] (-2.7,0) -- (2.9,0) node[right, font=\footnotesize] {$v_x$};
    \draw[->, thick] (0,-2.4) -- (0,2.9) node[above, font=\footnotesize] {$v_y$};
    \node[font=\scriptsize, anchor=north east] at (-0.05,-0.05) {$0$};
    % Recta de restriccion ORCA
    \draw[udcpink, thick] (-2.6,0.1) -- (2.4,2.6);
    \node[font=\scriptsize, text=udcpink, rotate=26] at (-1.35,0.95) {restricción ORCA};
    % Etiquetas de regiones
    \node[font=\scriptsize, text=udcpink!85!black] at (-1.3,2.1) {velocidades prohibidas};
    \node[font=\scriptsize, text=black!65] at (1.5,-1.2) {velocidades permitidas};
    % Velocidad preferida (cae en zona prohibida)
    \draw[->, very thick, draw=ficblue] (0,0) -- (0.9,2.2) node[above, font=\scriptsize, text=ficblue] {$v^{\mathrm{pref}}$};
    % Velocidad optima: proyeccion ortogonal de v_pref sobre la recta de restriccion
    \draw[dashed, thin, draw=black!55] (0.9,2.2) -- (1.04,1.92);
    \draw[->, very thick, draw=black] (0,0) -- (1.04,1.92);
    \fill[black] (1.04,1.92) circle (1.3pt);
    \node[font=\scriptsize, anchor=west] at (1.12,1.72) {$v^{\mathrm{new}}$};
  \end{tikzpicture}
  \caption[Resolución de la velocidad segura en el espacio de velocidades de ORCA.]
  {Resolución de la velocidad segura en el espacio de velocidades de ORCA.
  Cada agente vecino y la geometría inducen un semiplano de velocidades prohibidas (zona sombreada).
  Como la velocidad preferida $v^{\mathrm{pref}}$ cae en una región prohibida, el algoritmo escoge la velocidad factible $v^{\mathrm{new}}$ más próxima, sobre la frontera de la región segura.}
  \label{fig:orca}
\end{figure}

\subsection{Integración con el sistema de \textit{jobs}}

Por su naturaleza masivamente paralela (cada agente resuelve su propio problema), la simulación de evasión se integró sobre el \textit{Job System} y Burst.
Los agentes vecinos se localizan eficientemente mediante un \textbf{\textit{hash} espacial}, que evita comparar cada agente con todos los demás.
Además, los \textit{jobs} se planifican al final del fotograma y se completan al inicio del siguiente, de modo que su cómputo se solapa con el renderizado, minimizando su impacto sobre la tasa de fotogramas.

\section{Discusión global del sistema}
\label{sec:discusion-global}

Tras recorrer las distintas iteraciones, conviene contemplar el sistema en su conjunto.
El resultado es una canalización (\textit{pipeline}) claramente diferenciada en dos grandes momentos, ilustrada en la Figura~\ref{fig:flujo-datos}.

Por un lado, una fase de \textbf{preprocesado en el editor} (el \gls{bake}), que transforma la geometría de la escena en un \acrshort{svo} y lo serializa en disco.
Esta fase concentra deliberadamente el grueso del coste computacional (las consultas físicas, la construcción de la jerarquía y el enlazado de la vecindad), aprovechando que se ejecuta una sola vez y fuera del tiempo real.

Por otro lado, una fase de \textbf{ejecución en tiempo real}, que carga el volumen precalculado de forma inmediata y, sobre él, resuelve peticiones de ruta (búsqueda, atajo y suavizado) y simula la evasión local entre agentes.
El diseño de esta fase está dominado por la preocupación por el rendimiento: estructuras de datos compactas y contiguas, paralelización mediante el \textit{Job System} y Burst, y un uso cuidadoso de la memoria que minimiza la intervención del recolector de basura.

\begin{figure}[htp]
  \centering
  \resizebox{\textwidth}{!}{%
  \begin{tikzpicture}[font=\small, node distance=7mm and 9mm,
      contA/.style={rounded corners=5pt, draw=ficblue!70, very thick, fill=none, inner sep=8pt},
      contB/.style={rounded corners=5pt, draw=udcpink!70, very thick, fill=none, inner sep=8pt}]
    % --- Fase de preprocesado (editor): banda superior ---
    \node[paso] (escena) {Escena};
    \node[paso, right=of escena] (raster) {Rasterización};
    \node[paso, right=of raster] (svo) {\acrshort{svo}};
    \node[almacen, right=of svo] (disco) {Volumen\\en disco};
    % --- Fase de ejecucion (tiempo real): banda inferior ---
    \node[paso, below=16mm of escena] (carga) {Carga del\\volumen};
    \node[paso, right=of carga] (astar) {Búsqueda A*};
    \node[paso, right=of astar] (postp) {Post-\\procesado};
    \node[paso, right=of postp] (evasion) {Evasión\\local};
    \node[paso, right=of evasion] (movim) {Movimiento\\del agente};
    % --- Conexiones dentro de cada banda ---
    \draw[flecha] (escena) -- (raster);
    \draw[flecha] (raster) -- (svo);
    \draw[flecha] (svo) -- (disco);
    \draw[flecha] (carga) -- (astar);
    \draw[flecha] (astar) -- (postp);
    \draw[flecha] (postp) -- (evasion);
    \draw[flecha] (evasion) -- (movim);
    % --- Traspaso entre fases: el volumen serializado se vuelve a cargar ---
    \draw[flecha, rounded corners=3pt] (disco.south) -- ++(0,-5mm) -| (carga.north);
    % --- Contenedores de fase (solo borde, en capa normal) ---
    \node[contA, fit=(escena)(raster)(svo)(disco)] (cajaA) {};
    \node[contB, fit=(carga)(astar)(postp)(evasion)(movim)] (cajaB) {};
    \node[font=\footnotesize\itshape, text=ficblue, anchor=south west]
      at ([yshift=1.5mm]cajaA.north west) {Preprocesado (editor)};
    \node[font=\footnotesize\itshape, text=udcpink!85!black, anchor=north west]
      at ([yshift=-1.5mm]cajaB.south west) {Ejecución (tiempo real)};
  \end{tikzpicture}%
  }
  \caption[Flujo de datos global del sistema, del preprocesado a la ejecución.]
  {Flujo de datos global del sistema, del preprocesado a la ejecución.
  En el editor, la escena se rasteriza para construir el \acrshort{svo}, que se serializa en disco.
  En tiempo de ejecución, el volumen se carga y, sobre él, se resuelven las consultas de ruta (búsqueda, postprocesado y evasión local) que culminan en el movimiento del agente.}
  \label{fig:flujo-datos}
\end{figure}

A lo largo del desarrollo, varias decisiones de diseño transversales han demostrado su valor.
La representación dispersa y la codificación compacta de la información (las máscaras de bits de las hojas y los enlaces de 32 bits) reducen tanto el consumo de memoria como los fallos de caché.
La estrategia de lo grueso a lo fino, presente tanto en la rasterización como en la propia búsqueda de rutas, evita sistemáticamente trabajar a máxima resolución donde no es necesario.
Y la separación modular del código, establecida desde la primera iteración, ha permitido validar cada incremento de forma aislada, sosteniendo la fiabilidad del conjunto.
Esta separación, lejos de ser estática, creció de forma controlada a lo largo del proyecto: la Figura~\ref{fig:arquitectura-final} detalla la arquitectura final del módulo \textit{Runtime}, con sus cinco espacios de nombres definitivos —sus componentes principales y las dependencias que los relacionan— sobre el cimiento común que constituye \texttt{Core}.

\begin{figure}[htp]
  \centering
  \begin{tikzpicture}[
      font=\small,
      nsbox/.style={rounded corners=3pt, draw=udcpink!60!black, fill=white, thick,
                    align=center, text width=3.3cm, inner sep=5pt, minimum height=1.3cm,
                    font=\footnotesize},
      ancho/.style={text width=12.2cm},
      usedep/.style={-{Stealth[length=2.6mm]}, semithick, draw=ficblue!75!black}
    ]

    % --- Capa intermedia: construccion, busqueda y evasion ---
    \node[nsbox] (path) {\textbf{Pathfinding}\\[1pt]\scriptsize SVOPathfinder, PathSmoother,\\ SVORaycast (DDA), \textit{Api}};
    \node[nsbox, left=9mm of path] (builder)
      {\textbf{Builder}\\[1pt]\scriptsize SVOBuilder, SVORasterizer,\\ SVONeighborLinker,\\ \textit{Jobs}, \textit{Reporting}};
    \node[nsbox, right=9mm of path] (avoid)
      {\textbf{Avoidance}\\[1pt]\scriptsize AvoidanceSimulation,\\ \textit{Orca}, \textit{Spatial}};

    % --- Cimiento (Core) y capa superior (Unity), a lo ancho ---
    \node[nsbox, ancho, below=10mm of path] (core)
      {\textbf{Core}\\[1pt]\scriptsize MortonCode, SVO, SVONode, SVOLeaf, SVOLink, NeighborSet};
    \node[nsbox, ancho, above=10mm of path] (unity)
      {\textbf{Unity}\\[1pt]\scriptsize NavVolumeSpace, NavVolumeAgent, NavVolumeObstacle, \textit{Hashing}};

    % --- Dependencias (A -> B: A usa B) ---
    \draw[usedep] (builder.south) -- (builder.south |- core.north);
    \draw[usedep] (path.south)    -- (path.south    |- core.north);
    \draw[usedep] (avoid.south)   -- (avoid.south   |- core.north);
    \draw[usedep] (path.west)     -- (builder.east);
    \draw[usedep] (unity.south -| builder.north) -- (builder.north);
    \draw[usedep] (unity.south -| path.north)    -- (path.north);
    \draw[usedep] (unity.south -| avoid.north)   -- (avoid.north);

    % --- Contenedor Runtime ---
    \begin{scope}[on background layer]
      \node[contenedor, fit=(core)(builder)(path)(avoid)(unity),
            label={[font=\bfseries]above:Runtime}] (runtime) {};
    \end{scope}

    % --- Modulos Editor y Tests (dependen solo de Runtime) ---
    \node[modulo, fill=ficblue!15, draw=ficblue,
          above=14mm of runtime.north west, anchor=south west]
      (editor) {\textbf{Editor}\\[2pt]\footnotesize Herramientas de\\configuración y\\visualización};
    \node[modulo, fill=udcgray!40, draw=udcgray!80!black,
          above=14mm of runtime.north east, anchor=south east]
      (tests) {\textbf{Tests}\\[2pt]\footnotesize Pruebas \textit{EditMode}\\y \textit{PlayMode}};
    \draw[dep] (editor.south) -- (editor.south |- runtime.north);
    \draw[dep] (tests.south)  -- (tests.south  |- runtime.north);

  \end{tikzpicture}
  \caption[Arquitectura final del módulo \textit{Runtime} y sus dependencias internas.]
  {Arquitectura final del módulo \textit{Runtime} y sus dependencias internas.
  Sobre el cimiento común \texttt{Core} se asientan \texttt{Builder}, \texttt{Pathfinding} y \texttt{Avoidance} (este último incorporado en la última iteración), y \texttt{Unity} orquesta a todos ellos exponiendo el sistema al motor.
  Las flechas señalan las dependencias principales entre espacios de nombres; los módulos \textit{Editor} y \textit{Tests} siguen dependiendo únicamente de \textit{Runtime} (cf.\ Figura~\ref{fig:arquitectura}).}
  \label{fig:arquitectura-final}
\end{figure}

En definitiva, el sistema resultante satisface los objetivos planteados: representa de forma eficiente el espacio navegable tridimensional, calcula sobre él trayectorias fluidas y libres de colisiones, y lo hace a través de una herramienta integrada en el motor y accesible para el usuario final.
El análisis del rendimiento del sistema y las conclusiones derivadas se presentan en los capítulos siguientes.
